\section{Concrete parameters and every fixed integer power}
\label{sec:parameters}

\subsection{Recovering the balanced benchmark}

If all \(M\) intervals equal \([0,n]\), the arc minimum used here is
zero, so \(\gamma=0\) for this particular divisor. This does not
prove that the entire actual gcd has rate zero.
The root equation becomes
\[
Y=\tfrac12\cot\frac{\pi}{4M}.
\]
Substitution gives
\[
D_0=-M\log\left(2\sin\frac{\pi}{4M}\right)-m,\qquad
G_0=m+M\log2.
\]
With \(m=289\), the resulting exponent is
\(2028.145382310168\ldots\).
This is a useful check of the normalization and contour calculation.

\subsection{Two, seven, and twelve layers}

Three configurations illustrate the improvement:
\begin{center}
\begin{tabular}{lrrl}
\toprule
Configuration&\(m\)&\(L\)&Offsets and multiplicities\\
\midrule
Two layers&255&14&\(t=(0,1),\ \nu=(128,128)\)\\
Seven layers&199&39&\(t=(0,\ldots,6),\ \nu=(80,20,\ldots,20)\)\\
Twelve layers&186&69&\(t=(0,\ldots,11),\ \nu=(66,11,\ldots,11)\)\\
\bottomrule
\end{tabular}
\end{center}
Their rates for \(k=3\), rounded here only for readability, are:
\begin{center}
\begin{tabular}{lrrrr}
\toprule
Configuration&\(D_0\)&\(G_0\)&\(\gamma\)&\(E_3\)\\
\midrule
Two layers&955.116623&419.770987&27.777682&1762.608232\\
Seven layers&664.537005&322.700112&44.701426&1282.033606\\
Twelve layers&605.861295&301.032295&45.436001&1203.972439\\
\bottomrule
\end{tabular}
\end{center}
The double-layer result follows also from~\eqref{eq:twogamma}.
For the twelve-layer case the corrected rates are
\begin{align*}
D&=651.2972962706767693926036416599\ldots,\\
G&=255.5962938344585049378330527898\ldots,\\
D-2G&=140.1047086017597595169375360801\ldots.
\end{align*}
The strict conclusion is \(E_3<1203.973\), not a conclusion inferred
from the rounded entries of the table.

These examples arose from bounded searches and a limited multiplicity
refinement. They are not the solution of a globally optimized
variational problem. In particular, more layers are not known to be
necessary or optimal.

\subsection{The accepted L99 instance, step by step}
\label{subsec:l99}

Use \(L=99\), offsets \(t_j=j\), and base multiplicities
\[
v_0=16,\qquad v_j=1\ (1\le j\le27),\qquad
c_j=\frac{99-2j}{99}.
\]
The total base multiplicity and length are
\[
V=16+27=43,\qquad
S=\sum_{j=0}^{27}v_jc_j
 =16+27-\frac{2}{99}\frac{27\cdot28}{2}
 =\frac{389}{11}.
\]
At scale \(s=4\), this gives \(M=172\), \(m=171\), and total
relative interval length \(4S=1556/11\). There are twenty-eight
distinct interval types, rather than twenty-eight rows: repeated
types account for the much larger matrix dimension \(172\).

The formal proof uses the fixed rational height
\[
Y=\frac{1801}{20}=90.05.
\]
This is close to the maximizing height, but need not be the exact
root of \eqref{eq:rootpreview}. The contour proof is valid at any
fixed \(Y>1/2\); a lower bound on the decay at a convenient height
is enough. The one-sided certified estimates are
\begin{align*}
D_0(Y)&\ge\frac{9789173}{20000}=489.45865,\\
G_0&\le\frac{13452441}{50000}=269.04882,\\
\gamma&\ge\frac{2761813}{50000}=55.23626.
\end{align*}
Here the last line denotes an available divisor rate: for every
positive slack the actual divisor eventually has at least this
exponential rate minus the slack. Equality with the exact digamma
sum is not needed.

Adding the lower divisor rate to the decay bound and subtracting
it from the growth bound gives the usable rational pair
\[
D=\frac{54469491}{100000},\qquad
G=\frac{2672657}{12500}.
\]
These are deliberately conservative rates, not decimal truncations
asserted to equal the exact analytic optimum. Their final use can
be checked entirely with integer arithmetic:
\begin{align*}
D+G&=\frac{75850747}{100000},\\
D-2G&=\frac{54469491-2\cdot21381256}{100000}
      =\frac{11706979}{100000},\\
171\frac{D+G}{D-2G}
 &=\frac{171\cdot75850747}{11706979}
  =\frac{12970477737}{11706979}.
\end{align*}
For the stated upward rounding, cross multiplication gives
\[
110793\cdot11706979-100\cdot12970477737=3550647>0.
\]
This proves the strict comparison with \(110793/100=1107.93\)
without relying on a printed decimal expansion.

\paragraph{Where the divisor number comes from.}
For the base weights the accepted finite-window calculation gives
\(\gamma_0=2761813/200000=13.809065\). Its exact endpoint table
has 3004 open cells, of which 3000 carry a positive certified lower
cost. The prime-period prefix covers \(r=0,\ldots,31\); the finite
tail covers \(r=32,\ldots,999999\). Their lower bounds are
\[
\text{prefix}\ge\frac{1380525198339888}{10^{14}},\qquad
\text{tail}\ge\frac{381339545481}{10^{14}}.
\]
The sum is \(13.80906537885369\), leaving positive room above
\(13.809065\) for the interval-shrinking loss. Nonnegative periods
outside this finite window can be discarded. Scaling all weights
by four scales the certified divisor rate by four as well.
Section~\ref{sec:formalization} explains how the finite inequalities
are connected to the actual asymptotic divisor.

\paragraph{Why the Arb and Lean decimals differ slightly.}
The earlier Arb calculation at the optimizing height gives an
expression near \(1107.9267863\). The accepted exact rational
expression is \(1107.9269670681\ldots\). The small increase comes
from a rational trial height, finite divisor windows, and one-sided
analytic estimates. The theorem claims the latter number. A sharper
experimental decimal is not substituted into an accepted theorem.

\subsection{An explicit scale recipe for every fixed power}

For the twelve-layer shape let
\[
v_0=6,\quad v_j=1\ (1\le j\le11),\quad
c_j=1-\frac{2j}{69},\quad \nu_j=sv_j.
\]
Then
\[
M=17s,\quad m=17s-1,\quad
S=\sum_jv_jc_j=\frac{347}{23}.
\]
The arc partition is unchanged when all weights are multiplied by
\(s\), so
\[
\gamma=s\gamma_*,\qquad
\gamma_*=4.1305455731413661362664616341\ldots.
\]
Introduce the fixed constants
\begin{align}
g_*&=17+S\log2-\gamma_*,
 &g_*&=23.3269358031325913579848229808\ldots,\label{eq:gstar}\\
A_*&=\sum_jv_jc_j\log c_j,\nonumber\\
d_*&=S\log(S/\Theta)-A_*-17+\gamma_*.\label{eq:dstar}
\end{align}
Since \(f\le17\) for the base weights, the prime number theorem gives
\(\gamma_*\le17\). Thus \(g_*\ge S\log2>1\), proving
\(G=sg_*-1>0\) without requiring a decimal estimate.
The certificate more precisely verifies \(g_*>23\).

\begin{proposition}[A root-free admissibility recipe]\label{prop:allodd}
For any fixed integer \(k\ge2\), let
\begin{equation}\label{eq:srecipe}
s_k=\max\left\{1,\
\left\lceil\exp\left(
\frac{(k-1)g_*-d_*}{S}+1\right)\right\rceil\right\}.
\end{equation}
This choice satisfies \(D-(k-1)G>0\). More explicitly, put
\[
\underline D_s=s(S\log s+d_*)+1-\frac{\Theta^2}{12sS},
\qquad G_s=sg_*-1.
\]
Whenever \(\underline D_s>(k-1)G_s\), the following bound needs no
numerical root \(Y\):
\begin{equation}\label{eq:rootfree}
\mu(\pi^k)\le
(17s-1)\frac{\underline D_s+G_s}
 {\underline D_s-(k-1)G_s}.
\end{equation}
In particular this is valid at \(s=s_k\).
\end{proposition}
\begin{proof}
Evaluate the decay expression at the trial height \(Y_0=sS/\Theta\);
the optimized \(D\) is at least this trial value.
For every \(x\ge0\),
\(\arctan x\ge x-x^3/3\) and \(\log(1+x^2)\ge0\).
Consequently
\[
F_c(Y_0)\ge c\log(Y_0/c)+c-\frac{c^3}{12Y_0^2}.
\]
Since \(0<c_j\le1\), one has \(\sum_jv_jc_j^3\le S\).
The terms \(sS-\Theta Y_0\) cancel, giving
\[
D\ge s(S\log s+d_*)+1-\frac{\Theta^2}{12sS}
=\underline D_s.
\]
At \(s=s_k\),
\[
\underline D_s-(k-1)G_s
\ge sS+k-\frac{\Theta^2}{12sS}>0.
\]
Here \(S=347/23\), \(s\ge1\), and \(\Theta<2\), so the last
positivity is immediate.

For \(G_s\ge0\) the function
\(m(D+G_s)/(D-(k-1)G_s)\) decreases as \(D\) increases.
Replacing \(D\) by the lower bound \(\underline D_s\) therefore
gives~\eqref{eq:rootfree}.
\end{proof}

This proposition supplies a definite upper-bound recipe for
\emph{every fixed integer power \(k\ge2\)}, in addition to the optimized examples
in Corollary~\ref{cor:numerical}. It does not require an unbounded
parameter search to establish admissibility.

\subsection{Why the higher-power constants still grow rapidly}

Expanding the root equation for large \(s\) gives
\[
Y=\frac{sS}{\Theta}+O(s^{-1}),\quad
D=s(S\log s+d_*)+1+O(s^{-1}),\quad
G=sg_*-1.
\]
Thus the threshold for \(s\) grows exponentially with \(k\).
The explicit choice~\eqref{eq:srecipe} yields
\begin{equation}\label{eq:growth}
E_k=O(k\rho^k),\qquad
\rho=\exp(g_*/S)=4.69343992780744986545\ldots.
\end{equation}
To justify the order estimate directly, that choice has
\(s_k=O(\rho^k)\), a gap at least \(s_kS+O(1)\), and numerator
\((17s_k-1)(\underline D_{s_k}+G_{s_k})=O(ks_k^2)\).
All implied constants are fixed by this twelve-layer shape.
Equation~\eqref{eq:growth} describes the construction as \(k\)
varies; it must not be confused with the separate limit
\(n\to\infty\) used to prove each fixed-\(k\) bound.

For orientation, the certified optimized values before upward
rounding are
\[
\begin{array}{c|l}
k&E_k\\ \hline
3&1203.97243920634698117\ldots\\
5&42460.70466848845471678\ldots\\
7&1272102.82718073869325320\ldots\\
9&35398936.33162709355243184\ldots
\end{array}
\]
The number \(s\) was optimized only within one fixed shape.

\subsection{A simple all-power bound from the L99 shape}
\label{subsec:l99uniform}

The same root-free calculation can use the rational lower divisor
rate already accepted for the L99 base shape. Put
\[
V=43,\quad S=\frac{389}{11},\quad
\gamma_0=\frac{2761813}{200000},\quad
A=\sum_{j=0}^{27}v_jc_j\log c_j,
\]
and define
\[
g=V+S\log2-\gamma_0,\qquad
d=S\log(S/\Theta)-A-V+\gamma_0.
\]
For every positive integer scale \(s\), the available growth rate
and trial-height decay lower bound are
\begin{equation}\label{eq:l99rootfree}
G_s=sg-1,\qquad
\underline D_s=s(S\log s+d)+1-\frac{\Theta^2}{12sS}.
\end{equation}
To obtain these formulas, repeat the preceding proof with
\(V=43\) in place of \(17\), and use the lower divisor rate
\(s\gamma_0\) consistently on both sides. No new arc partition is
needed when only the common scale changes.

\begin{proposition}[A coarse uniform paper-level bound]\label{prop:uniformL99}
For every integer \(k\ge3\), the L99 shape with \(s=8^k\) gives
\[
\mu(\pi^k)<301\cdot8^k.
\]
This is a consequence of the written estimates and the available
scaled-family interface; it is not an accepted all-\(k\) Lean theorem.
\end{proposition}
\begin{proof}
The elementary inequalities
\(2/3<\log2<3/4\), \(35<S<36\), and
\(13<\gamma_0<43\) imply
\[
1<g<43+36(3/4)-13<60.
\]
Since \(c_j\le1\), one has \(A\le0\). Also
\(S/\Theta>e\), using \(\Theta<2\) and \(e<3\), and
\(S-43+\gamma_0>0\). Hence \(d>0\).

For \(s=8^k\), \(\log s=3k\log2\), so
\(S\log s>70k\). Write \(e_s=\Theta^2/(12sS)\), with
\(0<e_s<1\). Equation~\eqref{eq:l99rootfree} gives
\begin{align*}
\Delta_s&=\underline D_s-(k-1)G_s\\
 &=s\{S\log s+d-(k-1)g\}+k-e_s
 >10ks.
\end{align*}
Moreover \(0<G_s<60s\), and \(m=43s-1<43s\). Therefore
\[
m\frac{\underline D_s+G_s}{\Delta_s}
=m\left(1+\frac{kG_s}{\Delta_s}\right)
<43s(1+6)=301\cdot8^k.
\]
The trial height \(sS/\Theta\) is greater than \(1/2\), so the
contour estimate applies. The transfer theorem completes the proof.
\end{proof}

This simple formula avoids an open-ended search to establish the
all-power quantifier. It is intentionally loose: even at \(k=3\)
it is much larger than the accepted optimized value. Its exponential
growth is also worse than the classical \(B(k)=O(k\log k)\) in
Proposition~\ref{prop:NWpowers}. Its value is a transparent route
from one certified interval shape to a uniform statement, rather
than a new asymptotic record.

\subsection{Sensitivity of the bound and safe numerical substitution}

Fix \(m,k\) and write \(\Delta=D-(k-1)G>0\). Direct
differentiation gives
\[
\frac{\partial E_k}{\partial D}=-\frac{mkG}{\Delta^2},\qquad
\frac{\partial E_k}{\partial G}=\frac{mkD}{\Delta^2}.
\]
When \(G\ge0\), a lower bound for \(D\) and an upper bound for
\(G\) therefore give a safe upper bound for \(E_k\), provided the
conservative gap stays positive. Close to a zero gap the same
formulas show why a small rate error can produce a large exponent
error. Displaying many decimal places does not repair an
unjustified estimate on the wrong side.

If only the available divisor rate changes, with \(D_0,G_0\)
fixed, then
\[
\frac{\mathrm dE_k}{\mathrm d\gamma}
=-\frac{km(D_0+G_0)}{(D_0-(k-1)G_0+k\gamma)^2}.
\]
This makes the arithmetic gain quantitative. By contrast, adding
layers also changes \(m,D_0,G_0\), so its net effect has to be
recomputed. More divisibility alone is not an optimization criterion.

\subsection{What would be needed to reach 10--100?}

In the present transfer formula, \(G\ge0\) implies
\[
E_k=m\left(1+\frac{kG}{D-(k-1)G}\right)\ge m.
\]
The accepted L99 example has \(m=171\). Hence this particular
parameter choice cannot possibly give \(E_3<100\).
This is not an impossibility theorem for all nested shapes.

A more substantial change would keep the \emph{target polynomial
degree} fixed at three. If
\[
|P(\lambda)|\ge c_\varepsilon\Hgt(P)^{-\nu-\varepsilon}
\quad(P\in\ZZ[\ii][Y],\ \deg P\le3,\ P\ne0),
\]
then taking \(P(Y)=8qY^3+\ii p\) gives
\[
|P(\lambda)|=q|\pi^3-p/q|,\qquad \mu(\pi^3)\le\nu+1.
\]
Only this sparse binomial family is actually needed.
The thresholds \(\nu\le99\) and \(\nu\le9\) correspond to upper
bounds \(100\) and \(10\).
Type II Hermite--Pad\'e and multiple orthogonal polynomial
constructions are plausible research tools; examples of
bounded-degree logarithmic approximation appear in
Marcovecchio~\cite{Marcovecchio}. Their numerical constants at other
arguments cannot be transplanted to \(\log\ii\) without checking
the complex endpoint and its arithmetic.
